\documentclass[11pt,a4paper]{article}

\usepackage[utf8]{inputenc}
\usepackage[T1]{fontenc}
\usepackage[french]{babel}
\usepackage{lmodern}
\usepackage[margin=2.5cm]{geometry}
\usepackage{amsmath,amssymb,amsthm}
\usepackage{stmaryrd}
\usepackage{mathtools}
\usepackage{booktabs}
\usepackage{enumitem}
\usepackage{mydef}   % notations de l'article (Burman, Ern, Mottier)

\newtheorem{lemma}{Lemme}
\newtheorem{proposition}{Proposition}
\theoremstyle{definition}
\newtheorem{definition}{Définition}
\theoremstyle{remark}
\newtheorem{remark}{Remarque}

\newcommand{\patch}{\mathcal{Q}}
\newcommand{\patches}{\mathfrak{Q}_h}
\newcommand{\Kp}{\mathsf{K}}
\newcommand{\Fp}{\mathsf{F}}
\newcommand{\A}{\mathsf{A}}
\newcommand{\B}{\mathsf{b}}
\newcommand{\Sm}{\mathsf{S}}
\newcommand{\U}{\mathsf{U}}

\title{Condensation statique locale pour les méthodes HHO non conformes\\
stabilisées par extension polynomiale ou par agglomération}
\author{Note technique accompagnant l'application \texttt{unfitted\_HHO} de \textsc{diskpp}}
\date{}

\begin{document}
\maketitle

\begin{abstract}
On décrit l'élimination locale des inconnues de cellule (condensation statique) pour la méthode
HHO non conforme stabilisée par extension polynomiale de Burman, Ern et Mottier, avec les notations
de l'article. Contrairement au cas standard (et à la stabilisation par agglomération), les inconnues
de cellule d'une sous-cellule mal coupée interviennent dans les contributions locales de la cellule
qui la stabilise~: l'élimination ne peut pas se faire élément par élément, mais elle se fait
patch par patch, les patchs étant les composantes connexes de la relation $S\sim\N(S)$.
Le système condensé ne porte que sur les inconnues de face~; il est nettement mieux conditionné
que le système complet.
\end{abstract}

\section{Forme algébrique du problème discret}

On reprend les notations de l'article~: maillage $\mesh$, sous-cellules $\pmesh=\pmeshOK\cup\pmeshKO$,
opérateur d'appariement $\N:\TKOmesh\to\mesh$, espace $\Uhatz=\UTHone\times\UFhonez\times\UTHtwo\times\UFhtwoz$
et problème discret
\begin{equation}\label{eq:pb}
\text{trouver } \uhat\in\Uhatz \text{ tel que } a_h(\uhat,\what)=\ell_h(\what)\quad\forall\what\in\Uhatz,
\qquad a_h := b_h + s_h^\circ + s_h^\Gamma + s_h^\N .
\end{equation}
On sépare les inconnues de cellule $\U_\mathsf{C}:=\UTHone\times\UTHtwo$ des inconnues de face
$\U_\mathsf{F}:=\UFhonez\times\UFhtwoz$. Dans des bases de ces espaces, \eqref{eq:pb} s'écrit
\begin{equation}\label{eq:block}
\begin{pmatrix} \A_{\mathsf{CC}} & \A_{\mathsf{CF}} \\ \A_{\mathsf{FC}} & \A_{\mathsf{FF}} \end{pmatrix}
\begin{pmatrix} \mathsf{u}_\mathsf{C} \\ \mathsf{u}_\mathsf{F} \end{pmatrix}
= \begin{pmatrix} \B_\mathsf{C} \\ \B_\mathsf{F} \end{pmatrix},
\end{equation}
avec $\A$ symétrique définie positive puisque $a_h$ est symétrique et coercive sur $\Uhatz$
(section~4 de l'article). La condensation statique consiste à éliminer $\mathsf{u}_\mathsf{C}$~:
\begin{equation}\label{eq:schur}
\Sm\,\mathsf{u}_\mathsf{F}=\mathsf{g}, \qquad
\Sm:=\A_{\mathsf{FF}}-\A_{\mathsf{FC}}\A_{\mathsf{CC}}^{-1}\A_{\mathsf{CF}},\qquad
\mathsf{g}:=\B_\mathsf{F}-\A_{\mathsf{FC}}\A_{\mathsf{CC}}^{-1}\B_\mathsf{C},
\end{equation}
puis à reconstruire $\mathsf{u}_\mathsf{C}=\A_{\mathsf{CC}}^{-1}(\B_\mathsf{C}-\A_{\mathsf{CF}}\mathsf{u}_\mathsf{F})$.
Le complément de Schur $\Sm$ est symétrique défini positif. Tout l'intérêt est que $\A_{\mathsf{CC}}$ soit
diagonale par blocs de petite taille~: \eqref{eq:schur} se calcule alors localement, sans jamais
assembler~\eqref{eq:block}.

\section{Le cas standard~: élimination élément par élément}

Pour une méthode HHO conforme, $a_h=\sum_{T\in\mesh} a_T$ où $a_T$ ne fait intervenir que
$\uThat=(\uT,\ufT)$. Ainsi $\A_{\mathsf{CC}}=\operatorname{diag}(\A_{TT})_{T\in\mesh}$ et
\[
\Sm=\sum_{T\in\mesh} \mathsf{E}_T^\top\big(\A_{\partial T\partial T}-\A_{\partial T T}\A_{TT}^{-1}\A_{T\partial T}\big)\mathsf{E}_T ,
\]
où $\mathsf{E}_T$ extrait les inconnues de face de $T$~: chaque matrice locale est condensée avant
l'assemblage. Pour une cellule coupée $T\in\cutmesh$, les deux composantes $u_{T^1},u_{T^2}$ sont couplées
par le terme d'interface de $\GTo$ et par $s_h^\Gamma$~; on les élimine donc ensemble (bloc
$\A_{TT}$ de taille $2\dim\P^{k+1}$). Avec la stabilisation par agglomération, une cellule agglomérée
$\patch$ (réunion de mailles, sans inconnue sur ses faces intérieures) joue exactement le rôle de $T$~:
ses inconnues de cellule, des deux côtés de l'interface, ne sont couplées qu'à ses propres faces,
et la condensation est la condensation standard, cellule agglomérée par cellule agglomérée.

\section{Extension polynomiale~: élimination patch par patch}

\subsection{Contributions locales}

Regroupons les termes de $a_h$ par sous-cellule (c'est aussi l'organisation de l'implémentation).
Pour $(T,i)\in\pmeshOK$ et $\vhat,\what\in\Uhatz$,
\begin{align}
a^{\textup{OK}}_{T,i}(\vhat,\what) :={}& \kappa_i(\GTi(\vThatplus),\GTi(\wThatplus))_{T^i} \nonumber\\
&+ \kappa_i h_T^{-1}(\projT(\vTi)-\vdTi,\projT(\wTi)-\wdTi)_{\dTi} \nonumber\\
&+ \sumexi \eta_\N\kappa_i h_T^{-2}(\vSi-\vTi^+,\wSi-\wTi^+)_{T^i} \nonumber\\
&+ \delta_{i1}\,\kappa_1h_T^{-1}(\jumpvT_\Gamma,\jumpwT_\Gamma)_{\TG}, \label{eq:aOK}
\end{align}
et pour $(T,i)\in\pmeshKO$ (rappelons que $\GTi(\uThatplus)=\nabla\uTi$ dans ce cas),
\begin{align}
a^{\textup{KO}}_{T,i}(\vhat,\what) :={}& \kappa_i(\nabla\vTi,\nabla\wTi)_{T^i}
+ \kappa_i h_T^{-1}(\projT(\vTi)-\vdTi,\projT(\wTi)-\wdTi)_{\dTi} \nonumber\\
&+ \delta_{i1}\,\kappa_1h_T^{-1}(\jumpvT_\Gamma,\jumpwT_\Gamma)_{\TG}, \label{eq:aKO}
\end{align}
de sorte que
\[
a_h=\sum_{(T,i)\in\pmeshOK} a^{\textup{OK}}_{T,i}+\sum_{(T,i)\in\pmeshKO} a^{\textup{KO}}_{T,i}
\]
(le terme $s_h^\Gamma$ de chaque cellule coupée est compté une fois, du côté $i=1$~; dans
l'implémentation, $s_h^\circ$ est de plus multiplié par $k+1$). Le second membre $\ell_h$ se
décompose de la même façon.

\subsection{Le graphe de couplage des inconnues de cellule}

Lisons dans~\eqref{eq:aOK}--\eqref{eq:aKO} quelles inconnues de cellule chaque contribution couple.
\begin{itemize}[leftmargin=*]
\item $a^{\textup{KO}}_{T,i}$ ne fait intervenir que les inconnues de cellule de $T$
(les deux côtés, à cause du saut $\jumpvT_\Gamma$ si $i=1$).
\item $a^{\textup{OK}}_{T,i}$ fait intervenir les inconnues de cellule de $T$ (les deux côtés
si $T$ est coupée, par le terme $\delta_{i1}(\jumpuT_\Gamma,\q\SCAL\nG)_{\TG}$ de $\GTi$), mais aussi celles de
toutes les cellules $S\in\N_i^{-1}(T)$~: $u_{S^i}$ par $s_h^\N$ et par les termes
$(\udSi-\uSi,\q^+\SCAL\nS)_{\dSi}$ du gradient étendu, et $u_{S^1},u_{S^2}$ par le saut
$\jumpuS_\Gamma$ si $i=1$.
\end{itemize}
Les inconnues de cellule d'une cellule mal coupée $S$ apparaissent donc à la fois dans sa propre
contribution $a^{\textup{KO}}_{S,i}$ et dans celle, $a^{\textup{OK}}_{\N_i(S),i}$, de la cellule qui la
stabilise~: $\A_{\mathsf{CC}}$ n'est pas diagonale par cellule, et l'élimination élément par élément
n'est pas possible. Elle le devient si l'on regroupe les cellules couplées.

\begin{definition}[Patchs]
On note $\patches$ l'ensemble des classes d'équivalence de $\mesh$ pour la plus petite relation
d'équivalence telle que $S\sim\N(S)$ pour tout $S\in\TKOmesh$. Un élément $\patch\in\patches$ est
appelé un patch~; ses inconnues de cellule sont $(u_{T^1},u_{T^2})_{T\in\patch}$ et ses inconnues
de face celles de toutes les faces des cellules de $\patch$.
\end{definition}

Un patch est le plus souvent réduit à une cellule. Une cellule $T$ en image de $\N$ forme avec
$\N^{-1}(T)$ le patch $\{T\}\cup\N^{-1}(T)$, conformément à la remarque de l'article selon laquelle
\og all the cell unknowns attached to $T\cup\bigcup_{S\in\N^{-1}(T)}S$ are blocked together\fg.
La définition par composantes connexes est nécessaire car ces ensembles peuvent se chevaucher~:
avec l'étape~4 de l'appariement, $T:=\N_1(S)\in\TKOtwomesh$ entraîne $\N_2(T):=S$, et une cellule
peut être à la fois stabilisante d'un côté et stabilisée de l'autre. Les patchs se calculent en
temps quasi linéaire par union--recherche sur les paires $(S,\N(S))$.

\begin{lemma}[Localisation]\label{lem:loc}
Pour tout $(T,i)\in\pmesh$, toutes les inconnues de cellule qui interviennent dans la contribution
locale associée à $(T,i)$ appartiennent au patch contenant $T$. Par conséquent, avec les inconnues
de cellule numérotées patch par patch,
\[
\A_{\mathsf{CC}}=\operatorname{diag}\big(\A_{\patch\patch}\big)_{\patch\in\patches}.
\]
\end{lemma}

\begin{proof}
Les inconnues de cellule de $a^{\textup{KO}}_{T,i}$ sont celles de $T$. Celles de
$a^{\textup{OK}}_{T,i}$ sont celles de $T$ et des $S\in\N_i^{-1}(T)$, qui vérifient $S\sim T$. Deux
inconnues de cellule de patchs différents n'apparaissent donc jamais dans une même contribution,
et le coefficient correspondant de $\A$ est nul.
\end{proof}

\begin{lemma}[Inversibilité des blocs]
Chaque bloc $\A_{\patch\patch}$ est symétrique défini positif.
\end{lemma}

\begin{proof}
$\A_{\patch\patch}$ est la matrice de la restriction de $a_h$ au sous-espace de $\Uhatz$ formé des
fonctions dont les inconnues de face sont nulles et les inconnues de cellule sont nulles hors de
$\patch$. La restriction d'une forme bilinéaire symétrique coercive à un sous-espace reste
symétrique coercive.
\end{proof}

\subsection{Condensation et reconstruction}

Pour $\patch\in\patches$, notons $\Kp$ ses inconnues de cellule, $\Fp$ ses inconnues de face non
soumises à la condition de Dirichlet, et $\mathsf{D}$ les inconnues de face de Dirichlet, de valeurs
connues $\mathsf{u}_\mathsf{D}$. On assemble la matrice dense du patch en sommant les contributions
locales des sous-cellules de ses cellules,
\begin{gather*}
\begin{pmatrix} \A_{\Kp\Kp} & \A_{\Kp\Fp} \\ \A_{\Fp\Kp} & \A_{\Fp\Fp}\end{pmatrix}
:=\sum_{(T,i)\in\pmesh,\ T\in\patch}\big(\text{matrice de la contribution de }(T,i)\big),\\
\begin{pmatrix} \B_{\Kp} \\ \B_{\Fp}\end{pmatrix} := \sum_{(T,i)\in\pmesh,\ T\in\patch} \big(\text{second membre de la contribution de }(T,i)\big)
- \A_{\cdot\mathsf{D}}\,\mathsf{u}_\mathsf{D}.
\end{gather*}
D'après le lemme~\ref{lem:loc}, $\A_{\Kp\Kp}$ et $\B_\Kp$ sont exactement les blocs du système global
(aucune autre contribution n'y participe), tandis que $\A_{\Fp\Fp}$ et $\B_\Fp$ ne contiennent que la
part du patch. On calcule alors
\[
\mathsf{X}_\patch:=\A_{\Kp\Kp}^{-1}\A_{\Kp\Fp},\qquad \mathsf{y}_\patch:=\A_{\Kp\Kp}^{-1}\B_\Kp,
\qquad
\Sm_\patch:=\A_{\Fp\Fp}-\A_{\Fp\Kp}\mathsf{X}_\patch,\qquad \mathsf{g}_\patch:=\B_\Fp-\A_{\Fp\Kp}\mathsf{y}_\patch,
\]
et l'on assemble $\Sm=\sum_\patch \mathsf{E}_\patch^\top\Sm_\patch\mathsf{E}_\patch$,
$\mathsf{g}=\sum_\patch\mathsf{E}_\patch^\top\mathsf{g}_\patch$, ce qui donne exactement~\eqref{eq:schur}.
Après la résolution de $\Sm\,\mathsf{u}_\mathsf{F}=\mathsf{g}$, les inconnues de cellule se
reconstruisent patch par patch~:
\[
\mathsf{u}_\Kp=\mathsf{y}_\patch-\mathsf{X}_\patch\,\mathsf{E}_\patch\mathsf{u}_\mathsf{F}.
\]
Le système~\eqref{eq:block} n'est jamais formé. Le complément de Schur couple les faces de toutes les
cellules d'un même patch~: près de l'interface, son stencil est légèrement plus large que celui du
HHO standard, dans la même mesure que l'extension élargit celui de $\GTi$
(motif de creux de la figure~8 de l'article).

\begin{remark}[Agglomération]
Pour la stabilisation par agglomération, les patchs sont les cellules agglomérées et l'on retrouve
la condensation standard. Les deux condensations sont donc la même construction appliquée à des
structures de couplage différentes~; seule la définition des patchs change.
\end{remark}

\begin{remark}[Version algébrique]
Le même résultat s'obtient à partir de la matrice assemblée~\eqref{eq:block}~: les patchs sont alors
les composantes connexes du graphe de $\A_{\mathsf{CC}}$. Cette variante ne demande aucune
connaissance de l'appariement, mais elle assemble d'abord le système complet, qui compte environ
trois fois plus d'inconnues pour $k=3$.
\end{remark}

\section{Mise en \oe uvre et résultats}

Dans \texttt{unfitted\_HHO}, l'option \texttt{-c 1} (défaut) réalise la condensation locale~:
les patchs sont obtenus par union--recherche sur les paires $(S,\N(S))$ (une cellule agglomérée par
patch avec \texttt{-p1}), les contributions des sous-cellules d'un patch sont sommées dans une matrice
dense, les inconnues de Dirichlet sont passées au second membre, puis le patch est condensé par une
factorisation $\mathsf{LDL}^\top$. L'option \texttt{-c 2} réalise la variante algébrique et
\texttt{-c 0} résout le système complet.

\paragraph{Vérification.} Pour les solutions (56), (57) avec $\kappa_2=10^4\kappa_1$, et sur le cercle
et la fleur, $k=0,\dots,3$, niveaux $0$ à $3$, les erreurs obtenues avec \texttt{-c 0}, \texttt{-c 1}
et \texttt{-c 2} coïncident à $10^{-12}$ près en valeur absolue. Sur le cercle au niveau 3
($80\times80$ mailles), on obtient 6312 patchs. Le plus grand regroupe 90 inconnues de cellule pour
$k=3$, soit six sous-cellules, avec l'extension polynomiale, et 30 avec l'agglomération.

\paragraph{Conditionnement.} Le système condensé est beaucoup mieux conditionné que le système complet,
et presque insensible à la position de la coupe. Pour le cercle de rayon $R=1/3+2/32$ (configuration
la plus défavorable de la figure~12 de l'article) et le carré décalé de $\delta=0.5\cdot10^{-5}$, au
niveau~0, avec les bases orthonormées et $\eta_\N=20(k+1)^2$~:
\begin{center}
\begin{tabular}{lcccc}
\toprule
& $k=0$ & $k=1$ & $k=2$ & $k=3$ \\
\midrule
cercle, système complet & $4.5\cdot10^{2}$ & $1.5\cdot10^{3}$ & $1.1\cdot10^{4}$ & $1.1\cdot10^{5}$ \\
cercle, complet après mise à l'échelle diagonale & $1.7\cdot10^{2}$ & $9.9\cdot10^{2}$ & $3.5\cdot10^{3}$ & $1.6\cdot10^{4}$ \\
cercle, système condensé & $42$ & $1.2\cdot10^{2}$ & $1.8\cdot10^{2}$ & $2.7\cdot10^{2}$ \\
\midrule
carré, système complet & $5.2\cdot10^{2}$ & $1.8\cdot10^{3}$ & $5.0\cdot10^{3}$ & $1.2\cdot10^{4}$ \\
carré, système condensé & $45$ & $1.3\cdot10^{2}$ & $1.9\cdot10^{2}$ & $2.8\cdot10^{2}$ \\
\bottomrule
\end{tabular}
\end{center}

\paragraph{Coût.} Pour $k=3$ au niveau 4 ($160\times160$ mailles), le système condensé compte
$2.1\cdot10^{5}$ inconnues contre $5.9\cdot10^{5}$. Avec le même solveur (Pardiso), sa
résolution prend 1.07~s contre 2.35~s pour le système complet (et 4.6~s pour le système complet
résolu par \texttt{SimplicialLDLT}, comme dans ProtoN). La condensation locale ne coûte rien de
mesurable~: l'assemblage avec condensation (1.33~s) est même un peu plus rapide que l'assemblage du
système complet (1.41~s), qui manipule beaucoup plus de coefficients.

\end{document}
